\chapter*{Preface}
\addcontentsline{toc}{chapter}{Preface}

This book concerns a simple device and a less simple consequence. The device is a pair of glasses whose lenses open and close in step with a projector, admitting some frames and rejecting others. The consequence is that one screen, shown to one room, can carry several films at once, and that the people in the room need not see the same one.

The idea did not begin as a proposal about cinema in general. It began as a detail of a fictional production: a science-fiction thriller called \emph{Rarely Needed Protocols} and a satirical children's comedy called \emph{Barely Needed Protocols}, written to be screened simultaneously, sharing a complementary soundtrack, and separated only by the shutter timing of the glasses each viewer wore. The same notes proposed that the glasses could be mechanical, like the shutter of an old projector, and that which film a viewer saw would depend on which frame the shutter started on. A separate set of notes proposed the inverse arrangement: a gallery that appears empty until glasses with the correct blink sequence reveal the work, and the possibility of patenting the sequences so that only licensed glasses could see. Both sets of notes are archived in public repositories with their commit histories, and Appendix~G quotes them.

Those early notes already contained most of what this book develops: selection by timing rather than by content, the need for sound that serves more than one picture, the difference between choosing a version and being prevented from seeing one, and the uneasy fact that a mechanism for selection is also a mechanism for control. What they did not contain was an account of what kind of object a film becomes when it is designed from the start to be seen in several ways at once. Most of the book is an attempt to supply that account.

The argument proceeds from hardware to meaning. Part~I treats the screen as a carrier of interleaved streams and asks what the optics can and cannot do. Parts~II through~V treat the film as a family of related realizations and ask how such a family is written, scored, shot, and edited. Parts~VI through~IX turn to the audience and to the institutions that decide which realizations exist and who receives them. The book's central claim, stated in the Introduction and developed in \Cref{ch:synchronized-difference}, is that collective viewing does not require that everyone see the same thing. It requires that different experiences remain synchronized closely enough to count as one occasion.
